\subsection{Rigidity of the zero-mass double}
\label{subsec:static:zero-mass}

The metric $\gamma_*$ is complete and scalar flat, has a single
asymptotically Euclidean end of mass zero, and is $W^{2,s}_{\mathrm{loc}}$
for $n/2<s<n$.  Its possible nonsmoothness is confined to the joining
hypersurface and the filled point.  We will show that it is Euclidean.
The main step is to vary this metric without leaving the scalar-flat
class: nonnegative mass for both signs of a compact variation will
force its Ricci tensor to vanish.

The numerical Riemannian Penrose inequality in
Theorem~\ref{thm:riemannian-input} is the input.  It is supplied by
\emph{Conformal flow and the Riemannian Penrose inequality with
minimizing frontiers}.  We first construct the Laplace inverse needed
both to insert a Green-function pole and to correct scalar curvature.
Enclosing the pole by a minimizing frontier gives smooth nonnegative
mass.  We then smooth compact Sobolev defects and correct the scalar
curvature, while computing the change in mass.  Finally we remove the
weak regularity defects of the Ricci-flat double and use volume
comparison.  Only the numerical
inequality is needed from the Riemannian theorem.

Throughout this subsection $n\ge3$, $d=n-2$, $k=n-1$, and
$\omega=|S^{n-1}|$.  The Laplacian is $\Delta=\operatorname{div}\nabla$,
and the ADM mass has factor $1/(2k\omega)$.

\begin{theorem}[Riemannian numerical input {\cite[Theorem~1.1]{CompanionRiemannian}}]
\label{thm:riemannian-input}
Let $n\ge3$, $d=n-2$, $k=n-1$, and $\omega=|S^{n-1}|$.
Let $(E,h)$ be a connected complete \emph{enclosing-set exterior}:
the end-containing closed region outside a filled set, with its
compact full frontier $\Sigma$ included.  For $p,q\in E$, define the
possibly extended rectifiable-path distance
\[
 d_E(p,q)=\inf\{\operatorname{Length}_h(\gamma):
       \gamma\subset E\text{ is a rectifiable path joining }p\text{ to }q\},
 \qquad \inf\varnothing=+\infty.
\]
Completeness means that every finite-distance equivalence class
$\{q\in E:d_E(p,q)<\infty\}$ is complete for the restricted metric.
All points and components of $\Sigma$ remain in $E$, and their full
perimeter is counted.  This convention neither identifies $E$ with
the length completion of its open exterior nor asserts finite-length
access from the end to every singular frontier point.
Assume that the ambient metric $h$ is smooth on $E$ and through
$\Sigma$, and that the following conditions hold:
\begin{enumerate}
\item There is one Euclidean coordinate end with compact complement,
      and $h-\delta=O_2(r^{-\tau})$ for some $\tau>d/2$.
\item The scalar curvature satisfies $R_h\ge0$,
      $R_h\in L^1(E,dV_h)$, and $R_h=O(r^{-n-\beta})$ on the end
      for some $\beta>0$.
\item The full frontier $\Sigma$ is outer minimizing and locally
      perimeter minimizing as the boundary of the filled side.
      Its regular part is a smooth embedded minimal hypersurface,
      and its singular set has Hausdorff dimension at most $n-8$.
\end{enumerate}
Then, with $A$ the reduced-boundary perimeter of $\Sigma$, counting
every component,
\[
 m_{\rm ADM}(h)\ge \frac12\left(\frac{A}{\omega}\right)^{d/k}.
\]
The assertion includes $A=0$ and requires neither a spin assumption
nor connectedness of $\Sigma$.
\end{theorem}

Below we construct a nonempty detached
frontier in a smooth scalar-flat metric and verify these hypotheses
before applying the inequality.

\begin{lemma}[A compact-source Laplace inverse]
\label{lem:static:compact-source}
Let $\gamma$ be a complete locally uniformly elliptic
$W^{2,s}_{\mathrm{loc}}$ metric on a connected boundaryless manifold,
where $n/2<s<n$.  Suppose that it has compact complement to one
smooth end on which $\gamma-\delta=O_j(r^{-d})$ for every $j$.
Fix a compact set $K$ and $d/2<b<d$.  For every $f\in L^s$ supported
in $K$ there is a unique weak solution $v\in W^{2,s}_{\mathrm{loc}}$
of $\Delta_\gamma v=f$ tending to zero at the end.  It satisfies
\begin{equation}
 \|v\|_{W^{2,s}(K_1)}+
 \sup_{r\ge R}r^b\sum_{j=0}^2r^j
          \|D^jv\|_{L^\infty(\{r<|x|<2r\})}
       \le C\|f\|_{L^s(K)},
 \label{eq:static:inverse-estimate}
\end{equation}
where $K_1$ contains $K$ and the core.  Moreover
\begin{equation}
 v=A_f r^{-d}+O_j(r^{-d-\eta_f}),\qquad
 A_f=-\frac{1}{d\omega}\int_N f\,dV_\gamma,
 \label{eq:static:source-flux}
\end{equation}
for every fixed $0<\eta_f<\min\{1,b\}$.  The inverse bound is uniform
under sufficiently small $W^{2,s}$ perturbations supported in a fixed
compact set; the end is kept unchanged.
\end{lemma}

\begin{proof}
We describe the underlying global estimate to fix the solvability
class.  This manifold has the Euclidean Sobolev inequality
\begin{equation}
 \|\varphi\|_{L^{2n/(n-2)}}
       \le C\|d\varphi\|_{L^2},\qquad
       \varphi\in C_c^\infty(N).
 \label{eq:static:global-sobolev}
\end{equation}
To verify the compact part of this assertion, on the end write
\[
 \varphi(r,\theta)=-\int_r^\infty\partial_t\varphi(t,\theta)\,dt.
\]
Weighted Cauchy--Schwarz gives
\[
 |\varphi(r,\theta)|^2
 \le\frac{r^{-d}}d
       \int_r^\infty|\partial_t\varphi|^2t^{n-1}\,dt.
\]
Integration over a fixed end annulus bounds its $L^2$ norm by the
global Dirichlet integral.  A Poincar\'e inequality on a connected
compact region meeting this annulus gives the same bound on the
core.  Finite-chart Sobolev inequalities on the core, and the
Euclidean Sobolev inequality applied after an end cutoff, now prove
Equation~\eqref{eq:static:global-sobolev}.  Uniform metric
comparability suffices in this argument.

Let $\mathcal D^{1,2}$ be the completion of $C_c^\infty(N)$ in the
Dirichlet norm.  Since a compactly supported $L^s$ function belongs
to $L^{2n/(n+2)}$, Equation~\eqref{eq:static:global-sobolev} and the
Lax--Milgram theorem solve
\[
 \int_N\langle dv,d\varphi\rangle_\gamma\,dV_\gamma
       =-\int_N f\varphi\,dV_\gamma,
 \qquad v\in\mathcal D^{1,2}.
\]
They also bound $\|dv\|_2$ and $\|v\|_{2n/(n-2)}$ by $C\|f\|_s$.
Local Moser iteration for this divergence equation gives a local
$L^\infty$ bound, because $s>n/2$.  The local $W^{2,s}$ estimate
applies as follows.  In coordinates write
$\Delta_\gamma=a^{ij}\partial_{ij}+B^j\partial_j$.
Here $a$ is continuous and uniformly elliptic, and
$B\in L^p$ with $p=ns/(n-s)>n$.  On small balls, freeze $a$, use the
constant-coefficient Calder\'on--Zygmund estimate, and absorb its
oscillation.  The drift is absorbed by H\"older and interpolation,
\[
 \|B Dv\|_s\le\|B\|_p\|Dv\|_n,
 \qquad
 \|Dv\|_n\le\varepsilon\|D^2v\|_s+C_\varepsilon\|v\|_s,
\]
where the strict inequality $p>n$ permits interpolation below the
Sobolev exponent for $Dv$.  These scalar interior estimates can
first be applied to smooth metric and source approximations; their
uniform energy, Moser, and interior bounds pass to the weak solution.
This proves the core part of
Equation~\eqref{eq:static:inverse-estimate}; the scalar estimates used
here are those of~\cite{GilbargTrudinger2001}.  The compact bound controls
the inner boundary of the end.  We next turn it into the decay and
flux estimates that will determine the conformal mass change.

For sufficiently large $r$,
\[
 \Delta_\gamma r^{-b}
      =b(b-d)r^{-b-2}+O(r^{-b-d-2})<0.
\]
Compare $v$ and $-v$ with $Cr^{-b}$, choosing $C$ from the fixed
inner sphere bound.  The comparison is legitimate on the unbounded
end by testing with $(\pm v-Cr^{-b})_+$: this function has zero inner
trace and belongs to $\mathcal D^{1,2}$ since $2b>d$.
It follows that $|v|\le Cr^{-b}$.  Rescaled interior estimates for
the smooth homogeneous end equation give all derivative bounds
there.  A decaying homogeneous solution has neither a positive
maximum nor a negative minimum by the weak maximum principle, so
the inverse is unique in the stated class.

Finally,
$\Delta_\delta v=O_j(r^{-n-b})$ on the end.  Apply
Lemma~\ref{lem:static:poisson-tail} and integrate the weak equation
over large coordinate balls to obtain
\[
 -d\omega A_f
   =\lim_{r\to\infty}\int_{S_r}\partial_{\nu_\gamma}v\,dA_\gamma
   =\int_N f\,dV_\gamma.
\]
There is no interior boundary in this distributional identity.
All constants above stay bounded under small compact $W^{2,s}$
metric perturbations: their $C^0$ norms and ellipticity, H\"older
moduli of $a$, and $L^p$ norms of $B$ are uniformly controlled, and
the end is fixed.
\end{proof}

\begin{lemma}[Smooth nonnegativity from the numerical Penrose theorem]
\label{lem:static:smooth-nonnegative-mass}
Assume the numerical assertion of Theorem~\ref{thm:riemannian-input}.
Every complete smooth connected boundaryless scalar-flat manifold
with compact complement to one end satisfying
$\gamma-\delta=O_j(r^{-d})$ has nonnegative ADM mass.
\end{lemma}

\begin{proof}
Fix a smooth point $p$.  A local Dirichlet Green function, cut off
inside its coordinate ball, has distributional Laplacian
$-d\omega\delta_p+f_0$, where $f_0$ is smooth and supported away
from $p$.  Subtract the solution of
$\Delta_\gamma v=f_0$ supplied by
Lemma~\ref{lem:static:compact-source}.  The resulting function $G$
is harmonic off $p$, tends to zero at the end, and has the positive
pole $G=r_p^{-d}(1+o(1))$ in geodesic distance from $p$, with the
corresponding differentiated leading term.  The maximum principle
on the punctured manifold gives $G>0$.  Its distributional equation
and the same flux calculation as in
Equation~\eqref{eq:static:source-flux} give
\[
 \Delta_\gamma G=-d\omega\delta_p,
 \qquad G=r^{-d}+O_j(r^{-d-\eta_G})\quad\hbox{at infinity}.
\]
For $\varepsilon>0$ put
$\gamma_\varepsilon=(1+\varepsilon G)^{4/d}\gamma$ on
$N\setminus\{p\}$.  This is scalar flat and complete, with the
original end and a second complete end at the pole.  Near the pole,
the length factor is comparable to $r_p^{-2}$.

For fixed $\varepsilon$, sufficiently small spheres about $p$ have
strictly negative mean curvature toward the original end.  In fact,
writing $U=1+\varepsilon G$, the conformal mean-curvature law is
\[
 H_{\gamma_\varepsilon}
 =U^{-2/d}\left(H_\gamma+\frac{2k}{d}\partial_{\nu_\gamma}\log U\right),
\]
and the expression in parentheses is
$k/r_p-2k/r_p+o(r_p^{-1})<0$.
Choose one such sphere $S_\rho$.  Truncate the pole end at a
smaller sphere $S_{\rho_0}$, where $0<\rho_0<\rho$, retaining the
closed ambient exterior $M_{\rho_0}$.  This is a smooth scalar-flat
manifold with compact boundary and compact complement to the original
Euclidean end.  Its ordinary metric completeness follows from its
compact smooth core, smooth boundary charts, and complete
asymptotically Euclidean tail.  This ambient completeness is independent
of paths in the singular enclosure to be constructed.  We keep
$M_{\rho_0}$ as an ambient neighborhood on both sides of the minimizing
frontier that will be constructed.

For the enclosure construction itself, let $X_\rho$ be the original-end
exterior bounded by $B=S_\rho$, with $B$ included.  Apply the smooth
full-enclosure results of
\cite[Lemmas~2.1--2.3]{CompanionRiemannian} to
$(X_\rho,\gamma_\varepsilon)$.  These results assume a smooth connected
manifold complete with its nonempty compact smooth boundary included,
one Euclidean coordinate end with compact complement, and
$g-\delta=O_2(r^{-q})$ for some $q>0$.  A compact smooth filling $O$
is attached behind $B$, and perimeter is measured in the resulting
boundaryless manifold, counting all free-frontier components and any
contact with $B$.  In our application choose the metric on that
auxiliary filling to agree with $\gamma_\varepsilon$ in a collar
behind $S_\rho$, where the original metric is smooth.  Its extension
farther inside the filling has no role in the scalar-curvature
argument on the exterior.

The end condition holds because
$\gamma_\varepsilon-\delta=O_j(r^{-d})$; in particular, large
coordinate spheres have mean curvature $k/r+O(r^{-d-1})>0$.
The enclosure lemmas give a compact full-perimeter minimizer with
connected end exterior.  Since $H_B<0$ toward $X_\rho$, their
strict-barrier conclusion puts a fixed exterior collar of $B$ inside
the minimizing filled set.  Its frontier is therefore detached from
$B$, is outer minimizing, and is locally perimeter minimizing in the
smooth metric.  Its singular set has Hausdorff dimension at most
$n-8$.  These statements use only smooth obstacle geometry and full
perimeter, independently of original-data variation or equality.

Consequently this compact frontier is locally perimeter minimizing,
outer minimizing, and minimal, with singular set of Hausdorff
dimension at most $n-8$ by minimizing-boundary regularity
\cite{Simon2018}.  Its area is positive, and the smooth metric is
defined on a neighborhood of it.  Its original-end exterior has
nonnegative scalar curvature, integrable scalar curvature (in fact
zero), and the decay $d>d/2$ required by
Theorem~\ref{thm:riemannian-input}.  To give its enclosing-set ambient realization,
use the already constructed $M_{\rho_0}$: its compact boundary is
strictly behind the detached frontier and its metric is smooth through
that frontier.  Let $D$ denote the retained closed enclosing-set
exterior, including its entire frontier, and let $d_D$ be the possibly
extended path distance defined above, using $\gamma_\varepsilon$.
We verify completeness of every finite-distance class of $d_D$.

Fix such a class $C$ and a $d_D$-Cauchy sequence $(x_j)$ in $C$.
Choose a subsequence $(x_{j_\ell})$ with
$d_D(x_{j_\ell},x_{j_{\ell+1}})<2^{-\ell-1}$.  By the definition
of the infimum, these successive points can be joined by rectifiable
paths $c_\ell\subset D$ of length less than $2^{-\ell}$.
Concatenate the paths.  Their total length is finite and their tail
lengths tend to zero.  Ambient distances between points on a tail are
bounded by that tail's length.  Ordinary completeness of $M_{\rho_0}$
therefore gives an ambient endpoint $x_\infty$, and $x_\infty\in D$
because $D$ is closed.  Appending this endpoint gives a rectifiable
path in $D$ with length at most the sum of the constituent lengths.
The appended limit lies in the same finite-distance class: the
completed concatenation joins it to the fixed chain point $x_{j_1}$
with finite length.  Its tails give
\[
 d_D(x_{j_\ell},x_\infty)
 \le\sum_{r\ge\ell}\operatorname{Length}_{\gamma_\varepsilon}(c_r)
 \longrightarrow0.
\]
The Cauchy property and the triangle inequality then give
$d_D(x_j,x_\infty)\to0$ for the whole sequence.  Thus $C$ is complete.
No restriction of $D$ to the finite-distance class of its end has been
made; all frontier points and components, and their full perimeter,
remain included.  The argument asserts neither finite-length access
to an arbitrary singular frontier point nor an identification with
the length completion of the open exterior.  Thus the numerical assertion of
Theorem~\ref{thm:riemannian-input} implies
$m_{\rm ADM}(\gamma_\varepsilon)\ge0$.
The old-end expansion of $G$ and the ADM normalization give
\[
 m_{\rm ADM}(\gamma_\varepsilon)
      =m_{\rm ADM}(\gamma)+2\varepsilon.
\]
Letting $\varepsilon\downarrow0$ proves the claim.  No property of
the discarded pole end other than the displayed barriers is used.
\end{proof}

\begin{lemma}[Conformal correction and mass at compact Sobolev defects]
\label{lem:static:rough-nonnegative-mass}
Under the numerical assertion of Theorem~\ref{thm:riemannian-input}, let $\gamma$
satisfy the hypotheses of Lemma~\ref{lem:static:compact-source} and
be scalar flat in distributions.  Suppose that it is smooth outside
a compact set.  Then $m_{\rm ADM}(\gamma)\ge0$.
For a sufficiently small compact $W^{2,s}$ metric perturbation
$\gamma'$ there is a positive factor $U=1+v$, tending to one,
such that $U^{4/d}\gamma'$ is scalar flat.  Its mass is
\begin{equation}
 m_{\rm ADM}(U^{4/d}\gamma')
   =m_{\rm ADM}(\gamma')
     -\frac{1}{2k\omega}\int_N R_{\gamma'}U\,dV_{\gamma'}.
 \label{eq:static:conformal-mass}
\end{equation}
The factor, its end monopole, and this mass depend differentiably
on a differentiable $W^{2,s}$ curve of perturbations supported in a
fixed compact set.
\end{lemma}

\begin{proof}
Put $c_n=4k/d$.  The conformal equation is
\[
 -c_n\Delta_{\gamma'}(1+v)+R_{\gamma'}(1+v)=0.
\]
Use the inverse $T_\gamma$ of
Lemma~\ref{lem:static:compact-source} to write it as
\begin{equation}
 v=T_\gamma\left[
       c_n^{-1}R_{\gamma'}(1+v)
          -(\Delta_{\gamma'}-\Delta_\gamma)v\right].
 \label{eq:static:conformal-neumann}
\end{equation}
The perturbation is supported in a fixed compact set $K$ containing
the nonsmooth locus.  Because $R_\gamma=0$, both terms in brackets
are supported in $K$.  The source norm is $L^s$ and the solution norm is that in
Equation~\eqref{eq:static:inverse-estimate}.  These products are
bounded in the indicated spaces: a second-order coefficient
perturbation is small in $L^\infty$, drift perturbations are small
in $L^p$, $Dv\in L^p$, $p/2>s$, and $v\in L^\infty$.
Also $R_{\gamma'}$ is small in $L^s$ because
$\gamma'\to\gamma$ in $W^{2,s}$ and $R_\gamma=0$.
Thus Equation~\eqref{eq:static:conformal-neumann} is solved by a
Neumann series, with $\|v\|\to0$ as $\gamma'\to\gamma$.
In particular $U>0$, the corrected metric is complete, and its
scalar curvature vanishes distributionally.  The coefficient maps $\gamma\mapsto\gamma^{-1}$,
$\gamma\mapsto\Delta_\gamma$, and $\gamma\mapsto R_\gamma$
are differentiable between these spaces: their derivatives contain
second derivatives linearly and products of two $L^p$ first
derivatives, with $p/2>s$.  Differentiating the convergent operator
inverse in Equation~\eqref{eq:static:conformal-neumann} therefore
proves differentiability in the perturbation parameter.

On the common end $v$ is harmonic for $\gamma'$, so it has the
monopole expansion of
Equation~\eqref{eq:static:source-flux}.  The ADM calculation for
$U=1+A r^{-d}+o_1(r^{-d})$ gives a mass change $2A$, or equivalently
\[
 m_{\rm ADM}(U^{4/d}\gamma')-m_{\rm ADM}(\gamma')
  =-\frac{2}{d\omega}
       \lim_{r\to\infty}\int_{S_r}
                   \partial_{\nu_{\gamma'}}U\,dA_{\gamma'}.
\]
Integrating $\Delta_{\gamma'}U=c_n^{-1}R_{\gamma'}U$ proves
Equation~\eqref{eq:static:conformal-mass}.  The flux integral is
finite, and its compact-source expression also proves its
differentiability.

To prove nonnegativity for the original rough metric, smooth its
coefficients in finitely many charts covering its nonsmooth compact
set, with a partition of unity, leaving the end fixed.  The resulting
smooth positive definite metrics $\gamma_j$ converge in $W^{2,s}$
on the core and uniformly; their scalar curvatures tend to zero in
$L^s$.  The preceding construction provides $U_j\to1$ and smooth
complete scalar-flat metrics $U_j^{4/d}\gamma_j$.  Smoothness of
$U_j$ follows from its smooth scalar elliptic equation.  The end
expansion is $\delta+O_\ell(r^{-d})$ at every derivative order $\ell$, so
Lemma~\ref{lem:static:smooth-nonnegative-mass} applies.  Compact
smoothing leaves the uncorrected end mass unchanged, and
Equation~\eqref{eq:static:conformal-mass} gives
\[
 0\le m_{\rm ADM}(U_j^{4/d}\gamma_j)
          \longrightarrow m_{\rm ADM}(\gamma).
\]
This proves nonnegativity at the asserted regularity as well.
\end{proof}

We now have nonnegative mass at the precise regularity of the filled
double and under its compact conformally corrected perturbations.
Applying this statement to both signs of a variation converts its
zero mass into Ricci flatness before any smooth rigidity theorem is
used.

\begin{proposition}[Euclidean conformal double]
\label{prop:static:euclidean-double}
The filled metric $(N,\gamma_*)$ of
Lemma~\ref{lem:static:filled-double} is isometric to Euclidean
$\R^n$.
\end{proposition}

\begin{proof}
\smallskip\noindent
\textbf{Step 1: Mass variation and weak Ricci flatness.}
Let $v_{ij}$ be an arbitrary smooth symmetric tensor compactly
supported in a smooth patch of $\gamma_*$, and put
$\gamma_t=\gamma_*+tv$.  For both signs of sufficiently small $t$,
Lemma~\ref{lem:static:rough-nonnegative-mass} constructs a positive
factor $U_t$ such that $\widetilde\gamma_t=U_t^{4/d}\gamma_t$ is
scalar flat.  The same lemma gives
$m_{\rm ADM}(\widetilde\gamma_t)\ge0$.
Since $U_0=1$ and $m_{\rm ADM}(\gamma_*)=0$, this differentiable
mass has derivative zero at $t=0$.
Equation~\eqref{eq:static:conformal-mass} gives precisely
\begin{align*}
 0
 &= -\frac{1}{2k\omega}\int_N
          R'_{\gamma_*}[v]\,dV_{\gamma_*}\\
 &= \frac{1}{2k\omega}\int_N
          \langle\Ric_{\gamma_*},v\rangle_{\gamma_*}\,dV_{\gamma_*}.
\end{align*}
For the second equality use
$R'[v]=\divg\divg v-\Delta\tr v-\langle\Ric,v\rangle$;
the divergence terms integrate to zero because the test is
supported in that smooth patch.  Hence $\Ric_{\gamma_*}=0$ on every
smooth patch.  Its coordinate components belong to $L^s$, as proved
in Lemma~\ref{lem:static:filled-double}.  The seam and the filled
point have zero volume, so $\Ric_{\gamma_*}=0$ also as a
distribution everywhere.

\smallskip\noindent
\textbf{Step 2: Smoothness at the seam and the filled point.}
We include the regularity argument before applying a smooth
comparison theorem.  For a locally uniformly elliptic
$W^{2,s}$ metric with $n/2<s<n$, put $p=ns/(n-s)>n$.
It belongs to $W^{1,p}$, so local harmonic coordinates exist as
$C^1$ diffeomorphisms of class $W^{2,p}$ by
\cite[Corollary~2.8]{JulinLiimatainenSalo2017} with harmonic exponent
two.  In original coordinates their equation is
\[
 a^{ij}\partial_{ij}y^\ell+B^j\partial_jy^\ell=0,
 \qquad a\in W^{2,s},\quad B\in W^{1,s}\cap L^p.
\]
We justify the extra derivative without assuming strong regularity
of the differentiated coordinates.  Define the distributional operator
on $W^{1,p}_{\mathrm{loc}}$ functions by
\[
 Lz=\partial_i(a^{ij}\partial_jz)
           +(B^j-\partial_i a^{ij})\partial_jz .
\]
Its lower-order product belongs to $L^{p/2}\subset L^s$ locally;
on $W^{2,s}$ functions it equals
$a^{ij}\partial_{ij}z+B^j\partial_jz$.
For $z=\partial_k y^\ell\in W^{1,p}$, differentiating the scalar
coordinate equation by weak product rules therefore gives
\[
 Lz=F,\qquad
 F=-(\partial_k a^{ij})\partial_{ij}y^\ell
            -(\partial_kB^j)\partial_jy^\ell\in L^s.
\]
Indeed $Dy$ is bounded, $DB\in L^s$, and
$Da,D^2y\in L^p$ with $p/2>s$.

Choose a small coordinate ball $B_r$ and a cutoff
$\chi\in C_c^\infty(B_r)$ equal to one on a smaller ball.
Then $w=\chi z\in W^{1,p}_0(B_r)$ satisfies $Lw=H\in L^s(B_r)$:
the additional terms are
$2a^{ij}(\partial_i\chi)(\partial_jz)$ and
$(a^{ij}\partial_{ij}\chi+B^j\partial_j\chi)z$.
Construct $v\in W^{2,s}(B_r)\cap W^{1,s}_0(B_r)$ with $Lv=H$
by perturbing the constant-coefficient Dirichlet inverse with
coefficient matrix $a(x_0)$ at the ball center.  Its
second-derivative estimate and the scaled Sobolev inequality give
\[
 \|B Dv\|_{L^s(B_r)}
 \le C r^{\,1-n/p}\|B\|_{L^p(B_r)}
                  \|D^2v\|_{L^s(B_r)},\qquad
 1-\frac np=2-\frac ns>0 .
\]
Continuity of $a$ makes the other operator error at most
$C\|a-a(x_0)\|_\infty\|D^2v\|_s$.  Both errors are small
for sufficiently small $r$, so a Neumann series constructs $v$.

Sobolev embedding gives $v\in W^{1,p}$ with zero trace.  Thus
$q=w-v\in W^{1,p}_0(B_r)\subset H^1_0(B_r)$ satisfies
\[
 \partial_i(a^{ij}\partial_jq)
          +(B^j-\partial_i a^{ij})\partial_jq=0 .
\]
Testing with $q$, using ellipticity with constant $\lambda_0>0$
and the Euclidean Sobolev inequality, gives
\[
 \lambda_0\|Dq\|_2^2
 \le C\|B-\operatorname{div}a\|_{L^n(B_r)}\|Dq\|_2^2 .
\]
The coefficient norm is arbitrarily small on a small ball because
$B,Da\in L^p$ and $p>n$.  Hence $q=0$.  On the smaller ball
$z=v\in W^{2,s}$, proving $y\in W^{3,s}_{\mathrm{loc}}$.

The inverse coordinates have the same regularity: their third
derivatives contain $D^3y\in L^s$ and quadratic $D^2y$ terms in
$L^{p/2}\subset L^s$.  The metric transformation formula likewise
has highest terms $D^2\gamma_*$, $D^3(y^{-1})$, and products of
two $L^p$ factors.  Consequently the harmonic-coordinate metric
remains $W^{2,s}$.  To justify tensoriality, approximate the metric
in $W^{2,s}$ and the coordinate map in $W^{3,s}$.  The latter
convergence is also $C^1$, so the maps and their inverses remain
diffeomorphisms on smaller common charts.  Their transformed
metrics converge in $W^{2,s}$; the same product bounds imply
convergence of their Ricci tensors in $L^s$.  The smooth Ricci
transformation law therefore passes to the limit.  Composition
of the $L^s$ tensors causes no additional difficulty: approximate
the limiting tensor by a smooth tensor and use the uniformly
bounded Jacobians of the coordinate maps.

The harmonic-coordinate Ricci identity
\cite[Equation~(4.3)]{DeTurckKazdan1981} now gives the weak system
\[
 (\gamma_*)^{ab}\partial_{ab}(\gamma_*)_{ij}
       =Q_{ij}(\gamma_*^{-1},D\gamma_*),
\]
whose right side belongs to
$L^{s_1}$, where $s_1=ns/(2(n-s))>s$.
The principal coefficients are continuous and uniformly elliptic.
To obtain the higher regularity without assuming it in advance,
localize a component with a cutoff in a small harmonic-coordinate
ball.  Its equation has right side in $L^{s_1}$: besides the quadratic
term, the cutoff terms contain only the metric and its first
derivatives.  Solve that Dirichlet equation in $W^{2,s_1}$ by the
constant-coefficient inverse and the small oscillation of the
principal coefficients.  The difference from the original localized
component lies in $H^1_0$ and solves the homogeneous equation.
Writing the equation in divergence form introduces the drift
$-\partial_i(\gamma_*)^{ij}\in L^p$; its $L^n$ norm is small on
this ball.  The energy uniqueness argument just used for the
coordinate derivatives forces the difference to vanish.  Hence each
metric component belongs to $W^{2,s_1}$ on a smaller ball.
Repeating strictly improves the exponent until it exceeds $n$ in
finitely many steps.  Indeed while $s_j<n$ the reciprocal recurrence
is $n/s_{j+1}=2n/s_j-2$, starting below two.  If an exponent equals
$n$, the embedding into $W^{1,p'}$ for every finite $p'$ supplies
the next improvement above $n$.  Thus the metric becomes
$C^{1,\alpha}$; the displayed system and successive Schauder
estimates make it smooth.  The harmonic transition maps are then
smooth by their scalar harmonic equations.  This establishes a
smooth complete Ricci-flat Riemannian structure on the same metric
space, including the filled point and the seam.

\smallskip\noindent
\textbf{Step 3: Euclidean volume rigidity.}
The one asymptotically Euclidean end gives asymptotic volume ratio
one.  More explicitly, outside each sufficiently large fixed core
the metric lies between $(1-\epsilon)\delta$ and
$(1+\epsilon)\delta$, while the core contributes only a fixed
additive distance and finite volume.  Squeezing geodesic balls
between the corresponding coordinate balls and then sending
$\epsilon\downarrow0$ gives
\[
 \lim_{R\to\infty}
 \frac{\Vol_{\gamma_*}(B_{\gamma_*}(o,R))}{(\omega/n)R^n}=1.
\]
For a complete smooth manifold with nonnegative Ricci curvature,
Bishop--Gromov comparison makes this ratio nonincreasing; its
small-radius limit is also one.  The classical equality case of Bishop--Gromov comparison gives a
global Euclidean isometry; this volume-rigidity statement is recalled
in~\cite[p.~3237, before Theorem~1.1]{CavallettiManini}.
\end{proof}


